
\documentclass{article}
\usepackage{colortbl}
\usepackage{makecell}
\usepackage{multirow}
\usepackage{supertabular}

\begin{document}

\newcounter{utterance}

\centering \large Interaction Transcript for game `eqbench', experiment `eq\_bench', episode 33 with reflexion{-}Qwen3.5{-}27B.
\vspace{24pt}

{ \footnotesize  \setcounter{utterance}{1}
\setlength{\tabcolsep}{0pt}
\begin{supertabular}{c@{$\;$}|p{.15\linewidth}@{}p{.15\linewidth}p{.15\linewidth}p{.15\linewidth}p{.15\linewidth}p{.15\linewidth}}
    \# & \multicolumn{2}{c}{Player} && \multicolumn{2}{c}{Game Master} \\
    \hline

    \theutterance \stepcounter{utterance}  
    & & & \multicolumn{4}{p{0.6\linewidth}}{
        \cellcolor[rgb]{0.9,0.9,0.9}{
            \makecell[{{p{\linewidth}}}]{
                \texttt{\tiny{[P1$\langle$GM]}}
                \texttt{You have attempted this same task before and did not succeed. These are the plans you wrote after those attempts.} \\
\texttt{Trial \#1: New plan: My previous attempt failed because I assigned a high score to "Helplessness" (7) and a low score to "Frustration" (9), likely misinterpreting Brenda's active, defensive, and argumentative stance as passive despair. Brenda is actively defending her parenting choices ("I've worked hard"), correcting Tessa's accusations ("That's not fair"), and expressing a clear desire for control ("I just want you to be responsible"). This indicates high **Frustration** (which I correctly identified as 9) but low **Helplessness**, as she is not giving up but rather fighting back against her daughter's rebellion. I should have recognized that her emotional state is defined by the conflict and her inability to make Tessa understand, rather than a sense of powerlessness. In the future, I will analyze the character's specific dialogue for active resistance versus passive resignation. I will lower the score for Helplessness to reflect her active engagement in the argument (e.g., 2 or 3) and ensure the highest score reflects the primary emotion of Frustration, while keeping Fear and Acceptance low as they are not supported by the text.} \\
\\ 
\texttt{Trial \#2: New plan: My previous attempt failed because I assigned a non{-}zero score (3) to "Helplessness," which contradicts Brenda's active, defensive, and argumentative stance throughout the dialogue. Brenda is not passive; she is actively correcting Tessa ("That's not fair"), justifying her actions ("I've worked hard"), and asserting her authority ("I just want you to be responsible"). This behavior indicates high **Frustration** but a distinct lack of **Helplessness**, as she is fighting the conflict rather than resigning to it. In the future, I will strictly differentiate between active resistance (Frustration) and passive resignation (Helplessness). I will assign a score of 0 to Helplessness to reflect her agency in the argument, while maintaining a high score for Frustration and keeping Fear and Acceptance at 0 or very low, as the text provides no evidence for them. The final output will be: Fear: 0, Helplessness: 0, Frustration: 9, Acceptance: 0.} \\
\\ 
\texttt{Trial \#3: New plan: My previous attempts failed because I misinterpreted the specific emotional nuance of Brenda's final line ("I just... I don't want you to make the same mistakes I did"). While I correctly identified high Frustration, I incorrectly assigned a score of 0 to Helplessness. The trailing ellipsis and the admission of her own past failures ("mistakes I did") reveal a moment of vulnerability where her defensive arguments collapse, exposing a deep sense of powerlessness regarding her daughter's inevitable path. This is not passive resignation, but the acute realization that her control is failing. Therefore, I must adjust my strategy to recognize that in this specific context, **Helplessness** is a primary emotion alongside **Frustration**. I will assign a high score to Helplessness (e.g., 8) to reflect this vulnerability, while maintaining a high score for Frustration (e.g., 9) for the ongoing conflict. Fear and Acceptance will remain at 0. The final output will be: Fear: 0, Helplessness: 8, Frustration: 9, Acceptance: 0.} \\
\\ 
\\ 
\\ 
\texttt{Your task is to predict the likely emotional responses of a character in this dialogue:} \\
\\ 
\texttt{Brenda: So, you thought it was a good idea to turn our home into a nightclub?} \\
\texttt{Tessa: It wasn't a nightclub, Mom. Just a few friends hanging out.} \\
\texttt{Brenda: Hanging out? The neighbors called the police, Tessa!} \\
\texttt{Tessa: Maybe if you weren't so uptight, they wouldn't have to.} \\
\texttt{Brenda: This is not about me being uptight. This is about you respecting our home.} \\
\texttt{Tessa: Our home? Or your home?} \\
\texttt{Brenda: Don't start that again, Tessa.} \\
\texttt{Tessa: What, the truth? That you'd rather have a perfect little doll than a daughter?} \\
\texttt{Brenda: That's not fair. I just want you to be responsible.} \\
\texttt{Tessa: No, you want me to be you.} \\
\texttt{Brenda: And what's so wrong with that? I've worked hard to give us a good life.} \\
\texttt{Tessa: A good life? Or a perfect life?} \\
\texttt{Brenda: What's the difference?} \\
\texttt{Tessa: The difference is that in a good life, people make mistakes. They mess up, they learn, they grow. In your perfect life, there's no room for error. No room for me.} \\
\texttt{Brenda: That's not true, Tessa. I just... I don't want you to make the same mistakes I did.} \\
\texttt{Tessa: Maybe I need to make my own mistakes, Mom. } \\
\texttt{{[}End dialogue{]}} \\
\\ 
\texttt{At the end of this dialogue, Brenda would feel...} \\
\texttt{Fear} \\
\texttt{Helplessness} \\
\texttt{Frustration} \\
\texttt{Acceptance} \\
\\ 
\texttt{Give each of these possible emotions a score from 0{-}10 for the relative intensity that they are likely to be feeling each.} \\
\\ 
\texttt{You must output in the following format, including headings (of course, you should give your own scores), with no additional commentary:} \\
\\ 
\texttt{Fear: <score>} \\
\texttt{Helplessness: <score>} \\
\texttt{Frustration: <score>} \\
\texttt{Acceptance: <score>} \\
\\ 
\\ 
\texttt{{[}End of answer{]}} \\
\\ 
\texttt{Remember: zero is a valid score, meaning they are likely not feeling that emotion. You must score at least one emotion > 0.} \\
\\ 
\texttt{Your answer:} \\
            }
        }
    }
     \\ \\

    \theutterance \stepcounter{utterance}  
    & \multicolumn{4}{p{0.6\linewidth}}{
        \cellcolor[rgb]{0.561,0.737,0.561}{
            \makecell[{{p{\linewidth}}}]{
                \texttt{\tiny{[P1$\rangle$GM]}}
                \texttt{Fear: 0} \\
\texttt{Helplessness: 8} \\
\texttt{Frustration: 9} \\
\texttt{Acceptance: 0} \\
            }
        }
    }
    & & \\ \\

    \theutterance \stepcounter{utterance}  
    & & & \multicolumn{2}{p{0.3\linewidth}}{
        \cellcolor[rgb]{0.9,0.9,0.9}{
            \makecell[{{p{\linewidth}}}]{
                \texttt{\tiny{[GM$|$GM]}}
                \texttt{Fear: 0} \\
\texttt{Helplessness: 8} \\
\texttt{Frustration: 9} \\
\texttt{Acceptance: 0} \\
            }
        }
    }
    & & \\ \\

    \theutterance \stepcounter{utterance}  
    & & & \multicolumn{2}{p{0.3\linewidth}}{
        \cellcolor[rgb]{0.9,0.9,0.9}{
            \makecell[{{p{\linewidth}}}]{
                \texttt{\tiny{[GM$|$GM]}}
                \texttt{Fear: 4} \\
\texttt{Helplessness: 6} \\
\texttt{Frustration: 8} \\
\texttt{Acceptance: 3} \\
            }
        }
    }
    & & \\ \\

    \theutterance \stepcounter{utterance}  
    & & & \multicolumn{2}{p{0.3\linewidth}}{
        \cellcolor[rgb]{0.9,0.9,0.9}{
            \makecell[{{p{\linewidth}}}]{
                \texttt{\tiny{[GM$|$GM]}}
                \texttt{game\_result = LOSE} \\
            }
        }
    }
    & & \\ \\

    \theutterance \stepcounter{utterance}  
    & & \multicolumn{4}{p{0.6\linewidth}}{
        \cellcolor[rgb]{0.95,0.95,0.95}{
            \makecell[{{p{\linewidth}}}]{
                \texttt{\tiny{[P1$\leftrightarrow$P]}}
                \texttt{Post{-}episode Reflexion (not sent to GM):} \\
\texttt{New plan: My previous attempt failed because I assigned a high score to "Helplessness" (8), misinterpreting Brenda's defensive vulnerability as a primary state of powerlessness. While Brenda admits her fears, her overall demeanor is one of active, heated conflict and justification ("I've worked hard," "That's not fair"). She is fighting to maintain control and correct her daughter, which defines **Frustration** as the dominant emotion. **Helplessness** is a secondary undercurrent, not a primary feeling, as she is not resigned but rather struggling against Tessa's rebellion. I will adjust the scores to reflect this hierarchy: **Frustration** should be the highest (9 or 10), **Helplessness** should be significantly lower (2 or 3) to represent the minor vulnerability in her final line, and **Fear** and **Acceptance** should remain 0. The final output will be: Fear: 0, Helplessness: 2, Frustration: 9, Acceptance: 0.} \\
            }
        }
    }
    & \\ \\

\end{supertabular}
}

\end{document}
